% =============================================================================
\chapter{Literature Review}
\label{ch:literature}
% =============================================================================

Research on corporate disclosures examines whether information provided by
management helps investors understand a firm's performance and risk. This
chapter focuses on risk-factor disclosures and their relationship with
stock-return volatility. It reviews evidence on disclosure length,
risk-related vocabulary, and comparisons between firms' disclosures. These
studies provide the background for the three hypotheses developed in
Chapter~3.

\section{Risk disclosures and market risk}

Risk disclosures can inform investors about uncertain outcomes that are not
fully captured by accounting figures. However, providing more information does
not necessarily reduce perceived risk. A disclosure may explain an existing
uncertainty more clearly, but it may also reveal a risk that investors had
previously underestimated. This distinction helps explain why useful
disclosure can be associated with higher, rather than lower, subsequent
volatility.

\textcite{campbell2014} examine the information content of mandatory
risk-factor disclosures in 10-K filings. They find that firms facing greater
risk disclose more risk factors and that the types of risks discussed reflect
their underlying exposures. The disclosed information is also associated with
market-based risk measures, information asymmetry, and firm value. Their
findings suggest that risk-factor sections contain information about firms'
economic circumstances and cannot simply be dismissed as standard legal text.

\textcite{kravet2013} study changes in textual risk disclosures in annual 10-K
reports. Increases in disclosure are associated with higher stock-return
volatility and trading volume around and after the filing, as well as greater
dispersion in analysts' forecast revisions. Their evidence is consistent with
disclosures drawing attention to uncertainty rather than necessarily resolving
it. Unlike a comparison of disclosure levels across firms, their approach
examines annual changes in risk disclosure. These are related but different
questions: firms that consistently disclose more risks need not experience
higher volatility whenever they expand their disclosures.

\textcite{bao2014} use a topic model to identify different risk types in 10-K
disclosures. They find that many risk categories have no significant
relationship with investors' risk perceptions, while informative categories can
be associated with either increases or decreases. Their findings suggest that
an aggregate disclosure measure may combine information with different
implications.

Together, these studies support examining disclosure alongside subsequent
market risk, but they do not establish that more disclosure always has the
same meaning. They also require care in interpretation. A relationship between
disclosure and later volatility may reflect underlying business risk,
investors' responses to information, or both. Observing volatility after a
filing does not, on its own, identify the effect of the filing.

\section{Disclosure length and standard language}

Disclosure length is a simple measure of how much a company says about risk. It
is relatively easy to calculate and compare across filings. Its main
limitation is that it treats all words as contributing equally to disclosure
volume, regardless of whether they describe a distinct exposure, repeat earlier
statements, or provide general qualifications.

The findings of \textcite{campbell2014} provide a reason to expect a positive
relationship between the amount of risk discussion and subsequent risk.
However, this does not mean that a word count directly measures the number or
severity of a firm's risks. Length can also reflect reporting practices and the
way a company organizes its discussion.

\textcite{dyer2017} document changes in 10-K reports between 1996 and 2013.
Reports became longer and contained more boilerplate language, repetition, and
text retained from earlier years. Their topic analysis shows that disclosures
concerning fair value, internal controls, and risk factors account for much of
the increase in length. This evidence suggests that growth in disclosure volume
can reflect reporting requirements and accumulated text, rather than only
changes in economic risk. Their study concerns the full 10-K, so its findings
should not be interpreted as a direct test of Item~1A length and subsequent
volatility.

Legal considerations offer another explanation for lengthy and standardized
disclosures. \textcite{cazier2021} find that longer and more boilerplate risk
disclosures are less likely to be judged inadequate in judicial and regulatory
reviews. They also find that peer firms are more likely to borrow language
after courts have judged it adequate. This shows that shared wording can serve
a legal purpose. It does not establish whether that wording provides investors
with useful new information.

These studies suggest that length combines several influences. A long section
may reflect greater exposure to risk, more extensive explanation, or a cautious
reporting style. H1 therefore tests whether disclosure volume is associated
with subsequent volatility, without treating length as a direct measure of
disclosure quality. Examining both differences between firms and changes
within firms helps assess whether the association reflects persistent company
characteristics or changes in disclosure over time.

\section{Risk-word density and disclosure content}

Textual analysis offers ways to examine what companies write, rather than only
how much they write. A common approach counts words from predefined
dictionaries. The choice of dictionary matters because words can have
different meanings in financial reporting and everyday language.

\textcite{loughran2011} show that general-purpose dictionaries can misclassify
common financial terms. They develop word lists suited to financial text,
including uncertainty and litigious categories. The uncertainty list captures
imprecision and uncertainty, while the litigious list includes legal language
that does not always indicate an actual dispute. In their full-report analysis,
uncertainty-word proportions are positively associated with subsequent
volatility, while the result for litigious wording depends on how words are
weighted. This provides a reason to examine the two categories separately
rather than assume that they carry identical information.

A dictionary measure remains an approximation. Counting selected words does
not fully capture the meaning of the sentences in which they appear. The same
term may occur in a general warning or a detailed explanation of a specific
event. \textcite{loughran2016} emphasize the measurement and implementation
challenges involved in textual analysis. Their discussion supports
interpreting textual variables according to what they directly measure, rather
than assigning them broader meanings without validation.

This distinction is particularly important for disclosure specificity.
\textcite{hope2016} study how precisely firms describe disclosed risks. They
find that more specific risk-factor disclosures are associated with stronger
market reactions to 10-K filings and that analysts are better able to assess
fundamental risk. Their findings support the usefulness of detailed
disclosure, but they do not establish that a higher proportion of uncertainty
or litigious words measures greater specificity. Nor do they imply that useful
risk disclosure must be associated with lower subsequent volatility.

For the present thesis, risk-word density measures the share of Item~1A
vocabulary belonging to the selected dictionaries. It is not a direct measure
of clarity, specificity, or completeness. A higher density may indicate more
discussion of uncertain outcomes and legal exposures, but it may also reflect
standard cautionary wording. The latter possibility could weaken the
relationship with volatility; it does not by itself predict a negative
relationship.

Length and density must also be considered together. A section can contain
more risk-related words simply because it contains more words overall.
Dividing by section length changes the question from the amount of risk wording
to its share of the text. H2 examines whether this composition provides
additional information after controlling for length. Given the differences
between prior measures and the present specification, the hypothesis does not
impose a direction on the association.

\section{Peer comparisons and disclosure omissions}

A disclosure can also be evaluated relative to what other firms report.
Industry peers provide a possible benchmark because they may face similar
business conditions and exposures. Such a comparison asks whether a firm
covers vocabulary commonly used by comparable companies, rather than whether
its own disclosure is long or contains many dictionary words.

Peer comparisons already appear in the risk-disclosure literature.
\textcite{kravet2013} distinguish disclosure patterns shared within an industry
and year from those that differ at the firm level. Their results are less
pronounced for disclosures that deviate from the industry-year pattern. This
indicates that the relationship between a firm's language and its peers'
language can matter when interpreting disclosure measures. It does not,
however, establish that missing peer vocabulary represents withheld risk
information.

Selecting suitable peers is itself a measurement issue. \textcite{hoberg2016}
use product descriptions in 10-K filings to construct firm-specific competitor
groups that change over time. Their approach illustrates how textual
information can identify relationships that broad industry classifications may
miss. Their study concerns product markets, rather than omissions from
risk-factor sections, but it highlights a relevant limitation: firms assigned
to the same industry need not have identical businesses or risk exposures.

Comparisons of text can also be affected by document length. \textcite{brown2011},
studying annual changes in MD\&A, account for the mechanical relationship
between textual similarity and length. They find that disclosure revisions are
related to changes in firms' economic circumstances and to filing-date price
responses. Although their comparison is across years within a firm, rather than
across peers, it provides a useful methodological lesson: a text-difference
measure may partly reflect the amount of text available for comparison.

The omission measure in this thesis extends these ideas by examining the
weighted share of peer vocabulary absent from a firm's Item~1A. One possible
interpretation is that a larger gap reflects incomplete discussion of relevant
exposures. If so, it may be associated with greater uncertainty and higher
subsequent volatility. This is the proposed reasoning behind H3, rather than
an established finding of the studies reviewed above.

Several alternative explanations remain possible. A firm may not face the same
risks as its peers, may describe them using different words, or may have a
shorter section. Shared peer vocabulary may also reflect imitation rather than
shared exposure. The measure therefore captures a vocabulary gap, not verified
omissions of economic risks or deliberate silence. Testing its relationship
with volatility requires accounting for disclosure length and examining whether
the result is sensitive to how the gap is measured.

\section{Synthesis and link to the hypotheses}

The literature supports studying risk disclosures as a source of information
about firm risk, while showing that different textual features capture
different aspects of reporting. Disclosure volume can reflect underlying
exposure, but also reporting practices. Dictionary measures describe vocabulary
composition, while specificity concerns the detail with which risks are
explained. Peer comparisons provide another perspective, although differences
in wording do not necessarily indicate missing information.

This thesis brings three measures into a common regression framework. H1
examines the association between Item~1A length and subsequent volatility. H2
asks whether risk-word density adds information after length is accounted for,
including whether uncertainty and litigious wording behave differently. H3
examines whether a peer-relative vocabulary gap provides additional information
beyond both measures.

The study is therefore a focused comparison and extension of existing
approaches. It does not assume that any single measure captures overall
disclosure quality. Chapter~3 translates these questions into variable
definitions, formal hypotheses, and empirical specifications.

